\section{拡散モデルと生成モデル}

本章では，生成モデルを，観測データの分布あるいは条件付き分布を学習し，そこから新しいサンプルを生成する確率モデルとして統一的に捉える。画像生成や文生成の成功は広く知られているが，重要なのは，生成モデルが単なる出力装置ではなく，潜在変数を通じた表現学習，尤度に基づく学習原理，近似推論，逐次生成，さらには予測と計画へ接続される一般的な枠組みであるという点にある。

本章の目標は三つある。第一に，潜在変数モデルの学習という観点から，EMアルゴリズムと変分オートエンコーダ（variational autoencoder; VAE）を位置づける。第二に，ノイズ除去過程としての拡散モデルを，離散時間から連続時間へと見通しよく整理する。第三に，自己回帰モデルと世界モデルを通じて，生成モデルが系列生成や意思決定支援へどのように拡張されるかを理解する。以下では，「潜在変数モデルの学習」「逐次生成」「ノイズ除去過程」という三つの視点が一貫して現れるように議論を進める。

\subsection{生成モデルの考え方}

この概念が必要になるのは，分類や回帰だけでは扱えない問いが多く存在するからである。たとえば欠損補完，異常検知，反実仮想的なシミュレーション，都市活動の需要分布の生成，交通軌跡の合成などでは，入力に対してラベルを当てるだけでなく，データそのものがどのような確率法則に従って生じるかを記述する必要がある。生成モデルは，この要請に対して，観測 \(x\) の分布を直接あるいは間接に学習する立場をとる。

\subsubsection{分布を学ぶとは何か}

識別モデル（discriminative model）は，典型的には入力 \(x\) からラベル \(y\) への条件付き分布 \(p_{\theta}(y \mid x)\) を学習する。ここで \(\theta\) は学習すべきパラメータである。これに対して生成モデル（generative model）は，観測データ \(x\) の分布 \(p_{\theta}(x)\) を学ぶか，あるいは条件情報 \(c\) が与えられたときの条件付き分布 \(p_{\theta}(x \mid c)\) を学ぶ。記号 \(p_{\text{data}}(x)\) は真のデータ分布を表し，学習の目的は \(p_{\theta}(x)\) を \(p_{\text{data}}(x)\) に近づけることである。条件付き生成では，たとえば文脈 \(c\) に応じて文章や画像や将来軌跡を生成するために \(p_{\theta}(x \mid c)\) を考える。

生成モデルを学ぶ最も自然な原理は最大尤度推定（maximum likelihood estimation）である。独立同分布に従う観測データ \(x^{(1)},\dots,x^{(N)}\) が与えられたとき，尤度は
\begin{equation}
\prod_{n=1}^{N} p_{\theta}(x^{(n)})
\end{equation}
であり，対数尤度は
\begin{equation}
\sum_{n=1}^{N} \log p_{\theta}(x^{(n)})
\end{equation}
となる。ここで \(N\) は標本数である。対数をとると積が和に変わり，最適化や解析が行いやすくなる。したがって，生成モデルの中心課題は，対数尤度を大きくする \(\theta\) を求めることだといえる。

しかし，この方針は概念的には自然であっても，計算の上では難しい。高次元データに対して \(p_{\theta}(x)\) を直接評価すること自体が困難であり，とくに潜在変数を導入したモデルでは周辺化が障害となる。以下で見るように，EMアルゴリズム，VAE，拡散モデル，自己回帰モデルは，いずれもこの困難に対する異なる解法として理解できる。

\subsubsection{潜在変数モデル}

生成モデルで潜在変数（latent variable） \(z\) を導入する主な理由は，観測 \(x\) の背後にある低次元あるいは構造化された要因を分離し，表現学習を可能にするためである。潜在変数モデルの基本形は
\begin{equation}
p_{\theta}(x) = \int p_{\theta}(x \mid z)\, p(z)\, dz
\end{equation}
で与えられる。ここで \(p(z)\) は潜在変数 \(z\) の事前分布，\(p_{\theta}(x \mid z)\) は \(z\) が与えられたときの観測 \(x\) の条件付き分布であり，しばしばデコーダとも呼ばれる。積分は，潜在変数 \(z\) を明示的には観測しないまま，その全ての可能性を足し合わせる周辺化を表す。\(z\) が離散変数なら積分は総和に置き換わる。

この式は，観測データの複雑な分布を，潜在空間での比較的単純な分布と条件付き生成機構の組合せとして表現することを意味する。たとえば顔画像なら，姿勢，照明，表情といった要因を \(z\) に担わせることが期待される。交通軌跡なら，目的地，出発時刻帯，混雑回避傾向などを潜在表現として持たせることができる。潜在変数は観測されないため，その推論が必要になるが，その代わりにデータ分布の説明力と表現の可搬性が得られる。

問題は，周辺尤度 \(p_{\theta}(x)\) の計算が難しいことである。とくに \(p_{\theta}(x \mid z)\) をニューラルネットワークで表すと，積分は解析的に解けないことが多い。また，学習中に必要な事後分布 \(p_{\theta}(z \mid x)\) も容易には得られない。したがって，生成モデルにおける中心課題は，潜在変数を含む尤度最大化と近似推論をどのように両立させるかにある。

\subsubsection{本章の見取り図}

生成モデルには複数の流儀がある。敵対的生成ネットワーク（generative adversarial network; GAN）は識別器とのゲームを通じてサンプル品質を高める方法であり，自己回帰モデルは結合分布を条件付き分布の積に分解して逐次生成を行う。VAEは潜在変数モデルの学習を変分推論で実現し，拡散モデルはノイズ付加とノイズ除去の逐次過程として生成を定式化する。

本章では，まず潜在変数モデルを学習する古典的な方法としてEMアルゴリズムを扱い，次に厳密事後分布が扱えない深層モデルへ進むためにVAEを導入する。その後，ノイズ除去過程としての拡散モデルを離散時間と連続時間で整理し，最後に逐次生成の別系統として自己回帰モデルとLLMを概観する。さらに，生成モデルが表現学習と予測を通じて環境理解へ拡張される例として世界モデルを取り上げる。これにより，生成モデルを，単なるサンプル生成ではなく，潜在表現の獲得と確率的予測の一般原理として捉えられるようになる。

\subsection{EMアルゴリズム}

この概念が必要になるのは，潜在変数を含むモデルでは，尤度最大化がしばしば直接には解けないからである。観測データだけを見れば単純な最適化問題に見えても，実際には観測されていない変数を周辺化した対数尤度が現れ，対数と積分が絡み合うために計算が難しくなる。EMアルゴリズム（expectation-maximization algorithm）は，この困難を「潜在変数の推定」と「パラメータの更新」に分けることで扱う。

\subsubsection{完全データ尤度と不完全データ尤度}

潜在変数を \(z\)，観測を \(x\) とすると，完全データ（complete data）は \((x,z)\) であり，不完全データ（incomplete data）は \(x\) のみである。完全データの対数尤度
\begin{equation}
\log p(x,z \mid \theta)
\end{equation}
は，モデルによっては扱いやすい。一方，不完全データの対数尤度は
\begin{equation}
\log p(x \mid \theta) = \log \int p(x,z \mid \theta)\, dz
\end{equation}
であり，積分の外側に対数があるため，直接最適化しにくい。ここで \(\theta\) はモデルパラメータであり，\(p(x,z \mid \theta)\) は同時分布である。難しさの本質は，対数 \(\log\) が線形ではないため，潜在変数の周辺化と素直に交換できないことにある。

この問題を緩和するために，任意の分布 \(q(z)\) を導入して
\begin{equation}
\log p(x \mid \theta)
= \log \int q(z)\frac{p(x,z \mid \theta)}{q(z)}\, dz
\end{equation}
と書く。ここで \(q(z)\) は潜在変数 \(z\) に関する補助的な分布であり，後で適切に選ぶ。対数関数は凹関数であるから，Jensenの不等式より
\begin{equation}
\log p(x \mid \theta)
\ge
\int q(z)\log \frac{p(x,z \mid \theta)}{q(z)}\, dz
\end{equation}
が成り立つ。右辺は \(\theta\) に関する下界であり，しばしば自由エネルギー（free energy）あるいは変分下界と呼ばれる。

\subsubsection{EMの導出}

上の下界を
\begin{equation}
\mathcal{F}(q,\theta)
=
\mathbb{E}_{q(z)}[\log p(x,z \mid \theta)] + H(q)
\end{equation}
と書くことにする。ここで \(\mathbb{E}_{q(z)}[\cdot]\) は分布 \(q(z)\) に関する期待値であり，\(H(q)=-\mathbb{E}_{q(z)}[\log q(z)]\) はエントロピーである。この量は
\begin{equation}
\log p(x \mid \theta)
=
\mathcal{F}(q,\theta)
+
\mathrm{KL}\bigl(q(z)\,\|\,p(z \mid x,\theta)\bigr)
\end{equation}
と分解できる。右辺第二項の \(\mathrm{KL}(\cdot\|\cdot)\) はKullback--Leiblerダイバージェンスであり，二つの分布のずれを測る非負量である。したがって，\(\mathcal{F}(q,\theta)\) は常に対数尤度の下界であり，しかも \(q(z)=p(z \mid x,\theta)\) のときに下界は厳密に一致する。

EMアルゴリズムはこの事実を用いて，\(\theta\) と \(q\) を交互に更新する。現在のパラメータを \(\theta^{\text{old}}\) とすると，E-stepでは
\begin{equation}
q^{\text{new}}(z)=p(z \mid x,\theta^{\text{old}})
\end{equation}
と置く。これにより，下界は現在のパラメータに対して最もきつくなる。続いてM-stepでは，この \(q^{\text{new}}(z)\) を固定して \(\mathcal{F}(q^{\text{new}},\theta)\) を \(\theta\) について最大化する。エントロピー項 \(H(q^{\text{new}})\) は \(\theta\) に依らないため，M-stepは
\begin{equation}
Q(\theta \mid \theta^{\text{old}})
=
\mathbb{E}_{p(z \mid x,\theta^{\text{old}})}[\log p(x,z \mid \theta)]
\end{equation}
を最大化する問題になる。ここで \(Q(\theta \mid \theta^{\text{old}})\) は補助関数（auxiliary function）であり，古いパラメータで計算した事後分布の下で完全データ対数尤度の期待値をとったものである。

以上をアルゴリズムとしてまとめると，EMは次の二段階を反復する。
\begin{enumerate}
\item E-step: 現在のパラメータ \(\theta^{\text{old}}\) の下で事後分布 \(p(z \mid x,\theta^{\text{old}})\) を計算する。
\item M-step: 補助関数 \(Q(\theta \mid \theta^{\text{old}})\) を最大化して新しいパラメータ \(\theta^{\text{new}}\) を得る。
\end{enumerate}
この反復により，不完全データ対数尤度は単調非減少となる。EMが有効であるためには，E-stepで厳密な事後分布を計算でき，M-stepでも十分に扱いやすい最適化ができることが望ましい。

\subsubsection{GMMにおけるEMの直感}

混合ガウスモデル（Gaussian mixture model; GMM）はEMの典型例である。\(K\) 個のクラスタを考え，観測 \(x_n\) がクラスタ \(k\) に属する確率を混合係数 \(\pi_k\) で表す。潜在変数 \(z_n\) はデータ点 \(x_n\) のクラスタ割当を表し，\(z_n=k\) のとき
\begin{equation}
p(x_n \mid z_n=k,\theta)=\mathcal{N}(x_n;\mu_k,\Sigma_k)
\end{equation}
とする。ここで \(\mu_k\) はクラスタ \(k\) の平均，\(\Sigma_k\) は共分散行列である。

E-stepでは，各データ点が各クラスタに属する責務（responsibility）
\begin{equation}
\gamma_{nk}
=
p(z_n=k \mid x_n,\theta^{\text{old}})
=
\frac{\pi_k\,\mathcal{N}(x_n;\mu_k,\Sigma_k)}
{\sum_{j=1}^{K}\pi_j\,\mathcal{N}(x_n;\mu_j,\Sigma_j)}
\end{equation}
を計算する。責務 \(\gamma_{nk}\) は，データ点 \(x_n\) がクラスタ \(k\) にどれだけ説明されるかを表す重みである。M-stepでは，この重み付き平均として
\begin{equation}
N_k=\sum_{n=1}^{N}\gamma_{nk},\quad
\pi_k=\frac{N_k}{N},\quad
\mu_k=\frac{1}{N_k}\sum_{n=1}^{N}\gamma_{nk}x_n
\end{equation}
および
\begin{equation}
\Sigma_k
=
\frac{1}{N_k}\sum_{n=1}^{N}\gamma_{nk}(x_n-\mu_k)(x_n-\mu_k)^{\top}
\end{equation}
を更新する。ここで \(N_k\) はクラスタ \(k\) に割り当てられた有効サンプル数である。

この例から分かるように，EMは「見えない割当を現在のモデルで推定し，その推定結果に基づいてモデルを更新する」反復である。したがって，EMは潜在変数モデルの最も基本的な学習原理の一つといえる。ただし，深層生成モデルでは事後分布 \(p_{\theta}(z \mid x)\) を厳密に計算できないことが多い。その場合には，厳密事後分布を用いるEMの代わりに，近似事後分布を学習する変分推論が必要になる。次節では，この考え方をニューラルネットワークに組み込んだVAEを扱う。

\subsection{VAE}

この概念が必要になるのは，深層潜在変数モデルでは真の事後分布 \(p_{\theta}(z \mid x)\) がほとんどの場合に計算不可能だからである。EMは厳密な事後分布を扱えるときに強力だが，ニューラルネットワークを含む柔軟な生成モデルでは，E-stepそのものが閉形式で書けない。VAE\cite{Kingma2014,Rezende2014}は，近似事後分布を別のパラメータで学習することにより，生成モデルの尤度最大化を実用的な最適化問題へ変換する。

\subsubsection{変分推論とELBO}

VAEでは，真の事後分布 \(p_{\theta}(z \mid x)\) の代わりに，近似事後分布 \(q_{\phi}(z \mid x)\) を導入する。ここで \(\phi\) はエンコーダのパラメータであり，観測 \(x\) から潜在変数 \(z\) の分布を出力する。生成モデルは依然として
\begin{equation}
p_{\theta}(x,z)=p_{\theta}(x \mid z)p(z)
\end{equation}
と書かれ，\(p(z)\) は事前分布である。任意の \(q_{\phi}(z \mid x)\) に対して，
\begin{equation}
\log p_{\theta}(x)
=
\mathbb{E}_{q_{\phi}(z \mid x)}
\left[
\log \frac{p_{\theta}(x,z)}{q_{\phi}(z \mid x)}
\right]
+
\mathrm{KL}\bigl(q_{\phi}(z \mid x)\,\|\,p_{\theta}(z \mid x)\bigr)
\end{equation}
が成り立つ。右辺第二項は非負であるから，
\begin{equation}
\log p_{\theta}(x) \ge
\mathbb{E}_{q_{\phi}(z \mid x)}[\log p_{\theta}(x \mid z)]
- \mathrm{KL}(q_{\phi}(z \mid x)\|p(z))
\end{equation}
を得る。右辺がエビデンス下界（evidence lower bound; ELBO）である。ここで第一項は再構成項，第二項はKL正則化項である。

再構成項
\begin{equation}
\mathbb{E}_{q_{\phi}(z \mid x)}[\log p_{\theta}(x \mid z)]
\end{equation}
は，潜在変数 \(z\) から元の観測 \(x\) をどれだけよく説明できるかを測る。これはオートエンコーダにおける復元誤差に相当するが，確率的な形で定式化されている。一方，KL正則化項
\begin{equation}
\mathrm{KL}(q_{\phi}(z \mid x)\|p(z))
\end{equation}
は，データごとの近似事後分布を共通の事前分布 \(p(z)\) に近づける。これにより，潜在空間が無秩序に散らばることを防ぎ，サンプリング可能な構造を持たせる。VAEの学習は，この二つの要請，すなわち「よく再構成すること」と「潜在空間を整えること」の均衡である。

\subsubsection{再パラメータ化トリックと学習}

ELBOを確率的勾配法で最適化するには，\(q_{\phi}(z \mid x)\) からサンプルした \(z\) を通じて \(\phi\) に勾配を流す必要がある。しかし，単に「分布からサンプルする」と書くと，サンプリング操作が非微分的に見える。そこで再パラメータ化トリック（reparameterization trick）を用いる。

近似事後分布を平均 \(\mu_{\phi}(x)\) と標準偏差 \(\sigma_{\phi}(x)\) を持つ対角ガウス分布
\begin{equation}
q_{\phi}(z \mid x)=\mathcal{N}\bigl(z;\mu_{\phi}(x),\mathrm{diag}(\sigma_{\phi}(x)^2)\bigr)
\end{equation}
とすると，潜在変数 \(z\) を
\begin{equation}
z = \mu_{\phi}(x) + \sigma_{\phi}(x) \odot \epsilon,
\quad
\epsilon \sim \mathcal{N}(0,I)
\end{equation}
と書ける。ここで \(\odot\) は要素ごとの積，\(I\) は単位行列，\(\epsilon\) は標準正規分布からの補助ノイズである。重要なのは，確率性を \(\epsilon\) に押し込み，\(\mu_{\phi}(x)\) と \(\sigma_{\phi}(x)\) については決定論的な変換として \(z\) を表した点にある。これにより，\(\phi\) に関する勾配を，サンプル平均を通じて安定に推定できる。

事前分布を \(p(z)=\mathcal{N}(0,I)\) とし，近似事後分布を上の対角ガウスとした場合，KL項は解析的に計算できる。潜在次元を \(j=1,\dots,d\) とすると，
\begin{equation}
\mathrm{KL}(q_{\phi}(z \mid x)\|p(z))
=
\frac{1}{2}\sum_{j=1}^{d}
\left(
\mu_j(x)^2 + \sigma_j(x)^2 - \log \sigma_j(x)^2 - 1
\right)
\end{equation}
である。ここで \(\mu_j(x)\) と \(\sigma_j(x)\) は潜在変数の第 \(j\) 成分に対応する平均と標準偏差である。この閉形式は，VAEの実装を著しく簡潔にする。

\subsubsection{表現学習としてのVAEと限界}

VAEの重要な意義は，生成と表現学習を同時に行う点にある。潜在空間が事前分布により正則化されるため，近い潜在ベクトルはしばしば意味的に近いサンプルを生成する。したがって，潜在空間上の補間は滑らかな変化として解釈されやすく，表現の連続性が得られる。これは，画像や軌跡や都市形状のような高次元データを低次元の意味表現へ圧縮する上で有用である。

他方で，VAEには典型的な限界もある。第一に，デコーダが強すぎると近似事後分布 \(q_{\phi}(z \mid x)\) が事前分布 \(p(z)\) に近づきすぎ，潜在変数がほとんど情報を持たなくなることがある。これは事後崩壊（posterior collapse）と呼ばれる。第二に，尤度を重視した平均化的な学習により，再構成や生成結果がぼけやすい。とくに多峰的な分布を単純なガウス尤度で近似すると，複数の可能性を平均したような出力になりやすい。

それでもVAEは，EMの考え方を深層モデルへ拡張した重要な橋渡しである。EMが厳密事後分布を用いるのに対し，VAEは近似事後分布 \(q_{\phi}(z \mid x)\) 自体を学習する。この違いにより，潜在変数モデルの学習は大規模データと表現学習へ開かれた。次節では，潜在変数の推論ではなく，ノイズ除去の逐次過程として生成を定式化する拡散モデルを扱う。

\subsection{拡散モデル}

この概念が必要になるのは，高品質なサンプル生成と安定な学習を両立させる新しい生成原理が求められたからである。VAEは尤度と潜在表現を扱いやすい一方で出力がぼけやすく，GAN\cite{Goodfellow2014}は鋭いサンプルを生成できる一方で学習が不安定になりやすい。拡散モデル（diffusion model）\cite{Ho2020}は，データに徐々にノイズを加える前向き過程と，そのノイズを逆向きに除去する生成過程を組み合わせることで，尤度に根ざした学習と高品質生成を接続する。

\subsubsection{前向き過程と閉形式}

離散時間の拡散過程では，時刻 \(t=1,\dots,T\) に対して，観測 \(x_{t-1}\) に小さなガウス雑音を加えて \(x_t\) を得る。前向き過程は
\begin{equation}
q(x_t \mid x_{t-1})
=
\mathcal{N}(x_t; \sqrt{1-\beta_t}\,x_{t-1}, \beta_t I)
\end{equation}
で定義される。ここで \(\beta_t \in (0,1)\) は時刻 \(t\) のノイズ分散を表すスケジュールであり，\(I\) は単位行列である。\(\beta_t\) が小さいほど，1ステップ当たりの変化は小さい。しばしば \(\alpha_t = 1-\beta_t\) とおく。

この前向き過程を繰り返すと，元のデータ \(x_0\) は徐々に情報を失い，十分大きい \(T\) では標準正規分布に近い雑音 \(x_T\) へ到達する。重要なのは，累積係数
\begin{equation}
\bar{\alpha}_t = \prod_{s=1}^{t}\alpha_s
\end{equation}
を用いると，\(x_0\) から任意の時刻 \(t\) の状態 \(x_t\) を1ステップで直接サンプルできることである。具体的には
\begin{equation}
q(x_t \mid x_0)
=
\mathcal{N}(x_t; \sqrt{\bar{\alpha}_t}\,x_0, (1-\bar{\alpha}_t)I)
\end{equation}
となり，したがって
\begin{equation}
x_t=\sqrt{\bar{\alpha}_t}\,x_0+\sqrt{1-\bar{\alpha}_t}\,\epsilon,
\quad
\epsilon \sim \mathcal{N}(0,I)
\end{equation}
と書ける。ここで \(\epsilon\) は標準正規ノイズである。この閉形式は学習上きわめて重要であり，各時刻の中間状態 \(x_t\) を逐次シミュレーションせずに直接生成できる。

\subsubsection{逆過程とノイズ予測}

生成は前向き過程の逆，すなわち雑音 \(x_T\) から出発して徐々にノイズを除去し，最終的にデータ \(x_0\) を得る過程として定式化される。逆過程のモデルを
\begin{equation}
p_{\theta}(x_{t-1} \mid x_t)
\end{equation}
と書く。ここで \(\theta\) はニューラルネットワークのパラメータであり，この条件付き分布が真の逆向き遷移を近似するように学習する。前向き過程の各ステップが十分小さいとき，逆向き条件付き分布も近似的にガウスになるため，平均や分散をネットワークで予測する設計が自然になる。

元来，拡散モデルは変分下界に基づいて導出されるが，実用上はノイズそのものを予測する学習目標に簡約されることが多い。すなわち，前向き過程で \(x_0\) から \(x_t\) を生成する際に用いたノイズ \(\epsilon\) を，ネットワーク \(\epsilon_{\theta}(x_t,t)\) に予測させ，
\begin{equation}
\mathcal{L}_{\text{simple}}
=
\mathbb{E}\bigl[\|\epsilon - \epsilon_{\theta}(x_t,t)\|^2\bigr]
\end{equation}
を最小化する。ここで期待値は，データ \(x_0\)，時刻 \(t\)，およびノイズ \(\epsilon\) に関する平均である。損失が平均二乗誤差になるのは，真のノイズを回帰問題として学習する立場をとっているからである。

この三者の関係を明確にしておこう。前向き過程 \(q(x_t \mid x_{t-1})\) はデータを徐々に壊す既知の確率過程であり，逆過程 \(p_{\theta}(x_{t-1}\mid x_t)\) は壊れたデータから元の構造を回復する学習対象である。そしてノイズ予測 \(\epsilon_{\theta}(x_t,t)\) は，逆過程を実装するための実用的なパラメータ化である。すなわち，ネットワークが「いま見えている \(x_t\) の中にどれだけノイズが混ざっているか」を推定できれば，そこから一段前の状態 \(x_{t-1}\) を構成できる。

\subsubsection{条件付き生成と位置づけ}

条件付き生成では，条件 \(c\) を入力に加えて \(\epsilon_{\theta}(x_t,t,c)\) を学習し，\(p_{\theta}(x \mid c)\) に相当する生成を行う。画像生成ならテキスト条件，交通分野なら出発地，目的地，時刻帯，天候などを条件とできる。さらに，分類器なしガイダンス（classifier-free guidance）\cite{Ho2022}は，条件付き予測と無条件予測を組み合わせて条件への忠実度を調整する手法であり，生成品質と条件整合性の実用的な制御を可能にする。

拡散モデルの要点は，生成を一挙に行うのではなく，容易に記述できる破壊過程と，それを反転する学習可能な回復過程へ分解したことにある。この見方をさらに推し進めると，時間刻みを無限に細かくした連続時間の確率過程へ到達する。次節では，その連続時間化がどのように理論を統一するかを見る。

\subsection{連続時間拡散モデル}

この概念が必要になるのは，離散時間の拡散モデルをより一般的な確率過程として理解すると，理論とアルゴリズムの見通しが大きく改善するからである。離散時間のデノイジング拡散確率モデル（denoising diffusion probabilistic model; DDPM）では時刻 \(t=1,\dots,T\) の格子上でノイズを加えていたが，刻み幅を十分小さくすると，連続時間の確率微分方程式で前向き過程と逆過程を統一的に表現できる。

\subsubsection{前向きSDEと逆時間SDE}

連続時間では，時刻を \(t \in [0,T]\) とし，データ分布から出発する確率過程 \(x_t\) を確率微分方程式（stochastic differential equation; SDE）
\begin{equation}
dx = f(x,t)\,dt + g(t)\,d w
\end{equation}
で表すことが多い。ここで \(f(x,t)\) はドリフト項，\(g(t)\) は拡散係数，\(w\) は標準Wiener過程である。ドリフト項は決定論的な流れを，拡散項はランダムなノイズ注入を表す。前向きSDE（forward SDE）は，時刻とともに分布を単純化し，最終的に扱いやすい既知分布へ近づける過程に対応する。

このとき，各時刻における周辺分布を \(p_t(x)\) と書く。\(p_t(x)\) は時刻 \(t\) での \(x_t\) の確率密度であり，初期時刻 \(t=0\) ではデータ分布，終端時刻 \(t=T\) ではほぼガウス雑音になるように設計される。\citet{Song2021}が示した重要な結果によれば，適切な正則性の下で，時間を逆向きにたどる逆時間SDE（reverse-time SDE）は
\begin{equation}
dx=
\bigl[f(x,t)-g(t)^2\nabla_x \log p_t(x)\bigr]\,dt
+ g(t)\,d\bar{w}
\end{equation}
と書ける。ここで \(\nabla_x \log p_t(x)\) は時刻 \(t\) の分布の対数密度勾配であり，スコア関数（score function）と呼ばれる。また，\(\bar{w}\) は逆向き時間に対応するWiener過程である。つまり，逆過程を記述する鍵は，各時刻の密度そのものではなく，その勾配である。

\subsubsection{スコア関数とスコアベース生成}

スコア関数
\begin{equation}
\nabla_x \log p_t(x)
\end{equation}
は，空間の各点で「密度がどちらの方向に増えるか」を示すベクトル場である。スコアベース生成モデル（score-based generative modeling）\cite{Song2019}は，密度 \(p_t(x)\) を明示的にモデル化する代わりに，このスコア関数をニューラルネットワーク \(s_{\theta}(x,t)\) で近似する立場をとる。密度そのものを正規化定数まで含めて求めるのは難しいが，対数密度勾配は学習しやすい場合があり，しかも逆時間SDEの構成にはそれで十分である。

スコアマッチング（score matching）は，この勾配場を直接学習する枠組みである。とくにノイズを加えたデータに対するデノイジング・スコアマッチングでは，ノイズレベルごとのスコアを学習し，それを用いて逆向きにサンプリングする。離散時間の拡散モデルでノイズ \(\epsilon\) を予測する手法は，連続時間のスコア推定と密接に対応している。実際，ガウス摂動の下ではノイズ予測とスコア推定は定数倍の違いを除いて等価になる。

\subsubsection{離散時間との対応と統一的理解}

連続時間化の意義は，離散時間モデルと別物の理論を導入することではなく，それらを一つの枠組みへ包摂することにある。離散時間のDDPMで用いた \(\beta_t\) スケジュールは，連続時間では拡散係数 \(g(t)\) に対応し，逆過程の平均推定はスコア関数によるドリフト補正として理解できる。したがって，DDPMは連続時間SDEの時間離散化の一つとみなせる。

さらに，reverse-time SDEと同じ周辺分布を持つ常微分方程式
\begin{equation}
\frac{dx}{dt}
=
f(x,t)-\frac{1}{2}g(t)^2\nabla_x \log p_t(x)
\end{equation}
を考えることもでき，これは確率流常微分方程式（probability flow ODE）と呼ばれる。ここで常微分方程式はノイズ項を持たない決定論的な流れであり，尤度計算や高速サンプリングの解析に有用である。もっとも，本質は，連続時間化によって「ノイズ付加」「スコア推定」「逆過程サンプリング」が統一的に記述できるようになった点にある。

ここまでの拡散モデルは，ノイズ除去に基づく逐次生成であった。これに対して，系列の自然な順序に従って条件付き分布を直接学習する立場も存在する。次節では，言語モデルを代表例とする自己回帰生成を扱う。

\subsection{自己回帰モデルと LLM}

この概念が必要になるのは，系列データにはしばしば自然な因果順序があり，その順序に沿って結合分布を厳密に分解できるからである。文，音声，行動列，交通需要の時系列は，先行する要素が後続の要素に情報を与える。こうした場合には，潜在変数を積分消去するよりも，条件付き分布そのものを逐次的に直接学習する方が自然である。

\subsubsection{連鎖律による生成}

自己回帰モデル（autoregressive model）は，長さ \(T\) の系列 \(x_1,\dots,x_T\) の結合分布を確率の連鎖律により
\begin{equation}
p(x_1,\dots,x_T) = \prod_{t=1}^{T} p(x_t \mid x_{<t})
\end{equation}
と分解する。ここで \(x_{<t}\) は時刻 \(t\) より前の要素 \(x_1,\dots,x_{t-1}\) をまとめて表す。左辺の結合分布を直接扱う代わりに，各時刻の条件付き分布を学べばよいという点がこの枠組みの核心である。

学習は通常，教師ありの次トークン予測（next-token prediction）として行われる。すなわち，系列全体の対数尤度
\begin{equation}
\log p(x_{1:T})
=
\sum_{t=1}^{T}\log p(x_t \mid x_{<t})
\end{equation}
を最大化する。ここで \(x_{1:T}\) は系列全体を表す略記である。自己回帰モデルでは，各条件付き分布が明示的に規格化されているため，尤度計算が比較的直接的である。その代わり，生成時には \(x_1\) から順に1トークンずつサンプルしなければならず，本質的に逐次的である。

\subsubsection{n-gramからTransformerへ}

自己回帰言語モデルの古典的な形は \(n\)-gram であり，現在の単語の確率を直前 \(n-1\) 個の単語に条件づけて推定する。これは有限長文脈しか見られないという制約を持つが，連鎖律に基づく生成という発想自体はすでに備えている。次に再帰型ニューラルネットワーク（recurrent neural network; RNN）による言語モデルが現れ，隠れ状態に過去の情報を圧縮することで可変長文脈を扱えるようになった。

さらに，Transformer\cite{Vaswani2017}は自己注意機構により，系列内の離れた位置同士の依存関係を効率的に表現できるようにした。マスク付き自己注意を用いれば，将来のトークンを見ずに \(p(x_t \mid x_{<t})\) を学習できる。大規模言語モデル（large language model; LLM）は，この自己回帰Transformerを大規模データと大規模計算資源で学習したものであり，本質的には巨大な条件付き分布推定器である。

LLMにおいて，入力は通常，単語そのものではなく離散トークン列へ変換される。モデルはトークン埋め込みと位置情報を用いて文脈を表現し，次トークンの条件付き分布を出力する。パラメータ数とデータ量と計算量を増やすと性能が系統的に向上するというスケーリング則\cite{Kaplan2020}は，自己回帰モデルが高い表現力を持つことを示している。また，文脈内に与えられた例から振る舞いを調整するインコンテキスト学習（in-context learning）は，明示的なパラメータ更新なしに条件付き分布を柔軟に活用している現象として理解できる。

\subsubsection{拡散モデルとの対比}

自己回帰モデルと拡散モデルは，どちらも逐次生成を行うが，その逐次性の意味が異なる。自己回帰モデルでは，順方向の因果分解に基づいて \(x_1,x_2,\dots\) を一つずつ生成する。これに対して拡散モデルでは，ノイズの多い状態からノイズの少ない状態へと反復的に洗練する。前者は条件付き分布 \(p(x_t \mid x_{<t})\) を直接学ぶ枠組みであり，後者は逆過程 \(p_{\theta}(x_{t-1}\mid x_t)\) を学ぶ枠組みである。

この違いは，潜在変数や近似推論との関係にも表れる。自己回帰モデルは，潜在変数を明示せずに高次元分布を分解して扱う。一方，VAEや世界モデルでは潜在表現が前面に現れ，拡散モデルではノイズ状態という補助変数が逐次的に導入される。どの立場をとるかは，対象データの構造，必要とされる尤度計算，生成速度，条件付けの仕方に依存する。次節では，生成モデルがさらに一歩進み，環境の内部遷移そのものを学習する世界モデルへ向かう。

\subsection{世界モデル}

この概念が必要になるのは，生成モデルの役割が単なるサンプル生成にとどまらず，予測，計画，制御へ広がるからである。エージェントが環境中で行動するとき，本当に必要なのは「もっともらしい画像を生成すること」だけではない。観測の背後にある状態を圧縮して表現し，行動に応じて将来がどう変化するかを内部でシミュレーションできることが重要である。世界モデル（world model）は，この内部シミュレータを生成モデルとして構成する。

\subsubsection{潜在動的モデルの基本形}

時刻 \(t\) における観測を \(o_t\)，潜在状態を \(z_t\)，行動を \(a_t\) とする。世界モデルの基本構図は，観測列の背後に低次元の潜在状態列を仮定し，その時間発展を学習することである。典型的には
\begin{equation}
p(z_1)\,p(o_1 \mid z_1)\prod_{t=2}^{T} p(z_t \mid z_{t-1},a_{t-1})\,p(o_t \mid z_t)
\end{equation}
という形の潜在動的モデルを考える。ここで \(p(z_t \mid z_{t-1},a_{t-1})\) は遷移モデル，\(p(o_t \mid z_t)\) は観測モデルである。遷移モデルは行動 \(a_{t-1}\) を受けて潜在状態がどう変わるかを記述し，観測モデルは潜在状態からセンサ入力や画像や計測値がどう生成されるかを表す。

この構図は，VAEの潜在変数モデルを時間方向へ拡張したものと見なせる。観測 \(o_t\) をそのまま次時刻へ渡すのではなく，まず圧縮された潜在表現 \(z_t\) を介して扱うことで，高次元観測の冗長性を除きつつ，力学的に意味のある状態遷移を学習できる。表現学習と動的予測が結合している点が世界モデルの本質である。

\subsubsection{RSSMと想像上ロールアウト}

実際のモデルでは，確率的な潜在状態だけでなく，決定論的な再帰状態を併用することが多い。たとえば再帰状態空間モデル（recurrent state-space model; RSSM）は，確率的潜在変数とRNN的な隠れ状態を組み合わせ，部分観測環境でも十分な履歴情報を保持できるようにする。この種のモデルは潜在力学モデル（latent dynamics model）と総称され，モデルベース強化学習（model-based reinforcement learning）において重要な役割を果たしている。

学習された遷移モデルがあれば，実環境で行動せずとも，潜在空間内で将来軌跡を展開できる。これが想像上ロールアウト（imagined rollout）である。具体的には，現在の潜在状態推定 \(z_t\) から，候補行動列 \(a_t,a_{t+1},\dots\) を仮定して
\begin{equation}
z_{t+1}, z_{t+2}, \dots
\end{equation}
を遷移モデルで生成し，必要なら観測や報酬もデコードする。これにより，エージェントは現実世界で高価な試行錯誤を行う前に，内部モデル上で計画や方策評価を行える。交通シミュレーションでいえば，仮想的な需要変動や信号制御の変更が将来流動へ与える影響を，学習済み潜在力学上で検討することに相当する。

\subsubsection{生成モデルから意思決定支援へ}

世界モデルは，生成モデルの役割を三方向へ拡張する。第一に，観測の生成という意味で，依然として確率的なデコーダを持つ。第二に，潜在状態の学習という意味で，表現学習の装置でもある。第三に，行動条件付きの将来予測を通じて，計画と意思決定を支援する。ここでは生成が，見た目のサンプル合成ではなく，「もしこの行動をとったら何が起こるか」を内的に予測する能力へ転化している。

したがって，世界モデルは，本章で見てきた諸概念の自然な統合点である。潜在変数，近似推論，逐次生成，条件付き分布，表現学習がすべて現れるからである。生成モデルは，データ分布を学ぶ技術から出発しながら，最終的には環境の理解と行動選択を支える内部表現へと発展する。本章の締めくくりとして，この拡張可能性を意識しておくことが重要である。

\subsection*{まとめ}

本章では，生成モデルを，データ分布の学習という統一的な視点から整理した。まず，識別モデルとの対比を通じて，生成モデルが \(p_{\theta}(x)\) や \(p_{\theta}(x \mid c)\) を学び，新しいサンプルを生成するだけでなく，欠損補完や異常検知や合成データ生成に資する確率モデルであることを確認した。その上で，潜在変数モデル
\begin{equation}
p_{\theta}(x)=\int p_{\theta}(x \mid z)\,p(z)\,dz
\end{equation}
を導入し，複雑な観測分布を背後の表現と生成機構へ分解することの意味を見た。

続いて，潜在変数を含む尤度最大化の困難に対して，EMアルゴリズムが厳密事後分布を用いて下界最大化を行う方法であることを示し，さらにVAEが近似事後分布 \(q_{\phi}(z \mid x)\) を学習することで，深層潜在変数モデルへこの考え方を拡張することを見た。ここでは，尤度，近似推論，表現学習が密接に結びついている。次に，拡散モデルでは，生成がノイズ除去の逐次過程として再定式化され，離散時間のDDPMから連続時間のSDEとスコア関数へ至る統一的理解が得られた。さらに，自己回帰モデルでは，結合分布の因果分解に基づいて条件付き分布を逐次的に直接学習する立場を確認し，LLMがその大規模な実例であることを位置づけた。最後に，世界モデルを通じて，生成モデルが観測生成だけでなく，将来予測，想像上ロールアウト，計画へ拡張されることを見た。

以上を通じて，生成モデルは，データ分布の学習，潜在表現の獲得，逐次的・確率的生成，さらには予測と計画へ拡張される統一的枠組みとして理解できる。個々の手法は，潜在変数をどう導入するか，尤度や下界をどう最適化するか，生成をどの方向の逐次過程として捉えるかという点で異なるが，いずれも「データがどのように生じるか」を数理的に記述しようとする試みである。
